\documentclass[10pt,a4paper]{amsart}
\usepackage[margin=19mm]{geometry}
\usepackage{amsmath,amssymb,mathtools}
\usepackage[scr=rsfs,cal=euler]{mathalfa}
\usepackage{xfrac}
\usepackage[unicode,hidelinks,bookmarks=false]{hyperref}
\newcommand{\ie}{i.e.\ }
\newcommand{\cf}{cf.\ }
\newcommand{\resp}{respectively\ }
\newcommand{\Subsection}{\subsection}
\providecommand{\nobreakdash}{\nobreak\hbox{-}}
\setlength{\emergencystretch}{2em}
\multlinegap0pt
\usepackage{bm}
\usepackage{bbm}
\usepackage{dsfont}
\usepackage{xspace}
\usepackage{cancel}
\usepackage{mathtools}
\usepackage{amsthm}
\makeatletter
\@ifundefined{theorem}{\newtheorem{theorem}{Theorem}}{}
\makeatother
\title[On balanced biregular cages]{On balanced biregular cages}
\author{Araujo-Pardo Gabriela, Kiss Gy\"orgy}
\date{Source: arXiv:2604.22395v2 (CC BY 4.0)}
\begin{document}
\maketitle

\noindent\emph{Source and licence.} On balanced biregular cages. Version arXiv:2604.22395v2 (CC BY 4.0), distributed under the Creative Commons Attribution 4.0 International licence, as recorded in the arXiv licence metadata for this exact version and shown on the abstract page https://arxiv.org/abs/2604.22395v2. Full rights notice: licenses/arxiv-2604.22395-CC-BY-4.0.txt.\par\smallskip

\noindent\textit{Editorial scope.} Counted unit: the Theorem whose source label is \texttt{equal-5-6} in the article (source0.tex lines 373--381, source label \texttt{equal-5-6}), delivered together with the article's own complete original proof of it (source0.tex lines 383--443, one \texttt{proof} environment of 61 lines). No mathematics is written, repaired or re-derived: statement and proof are the article's own text line for line. Required context retained uncounted: lowerboundbabigraph (lines 291--312); tree-5 (lines 353--356). Every definition, display and label of those lines is retained unchanged. Omitted from this package and not redistributed: 1--290, 313--352, 357--372, 382--382, 444--1693. Nothing inside a retained excerpt is omitted or altered: every retained body is delivered exactly as the article's lines are written, so no omission receipt of any kind accompanies this package. Citations: the retained excerpts carry no citation command at all, so no bibliography entry of the article is transcribed and no reference is redistributed. Reference adaptation: none was needed and none was made; every reference inside the delivered unit resolves inside it, so no unresolved reference or citation is produced. Printed slips kept literal: no printed word, symbol, hypothesis or number of the counted statement or of its proof was altered. Layout and class: the article's own class is \texttt{article} with packages including t1enc, inputenc, amssymb, amsmath, amsthm, graphicx, hyperref, pgf, tikz, color, mathrsfs; the wrapper is the lane's pinned \texttt{amsart} editorial wrapper at 10pt on A4, which restates the theorem-like environments the retained text needs and adds the article's own macro definitions that the retained text needs (none), with extra wrapper packages (bm, bbm, dsfont, xspace, cancel, mathtools). The delivered numbering is the unit's own. Assets: no retained span contains a figure environment, a graphics-inclusion command, a PDF-page inclusion, a table or a TikZ picture; no image file is copied into this package, the package declares no asset dependency and no image file is redistributed with it. Licence: version arXiv:2604.22395v2 of this article is distributed under the Creative Commons Attribution 4.0 International licence, as recorded in the arXiv licence metadata for this exact version and shown on the abstract page https://arxiv.org/abs/2604.22395v2. A full rights notice is delivered at licenses/arxiv-2604.22395-CC-BY-4.0.txt. Characters: every retained excerpt is pure ASCII, so the delivered engine, XeLaTeX with its default Latin Modern text font, reports no missing character.
\begin{theorem}\label{lowerboundbabigraph}
Let $2\leq r<s$ be integers.
Then the order of an $(r,s;g)$-babi-graph is at least:
%$n_{bb}(r,s;g)$, where
   % $$n_{bb}(r,s;g)\geq
\begin{itemize}
\item
$1+s, \text{ if } g =3,$
\item
$2s, \text{ if } g =4,$
\item
$1+s + 
(rs+(s-r)^2-s)(1+(r-1)+\ldots + (r-1)^{\frac{g-5}{2}}),
\text{ if } g>3 \text{ odd}$,  
\item
$2\left( s+
(rs+(s-r)^2-2s+1)(1+(r-1)+\ldots + (r-1)^{\frac{g-6}{2}})\right)  ,
\text{ if } g>4 \text{ even.}$
\end{itemize}

%If equality holds, then $s=r+1$,  $f=\frac{v}{2}$ and the fat edges give a matching of %the fat vertices.
\end{theorem}

\begin{equation}
\label{tree-5}
|G_2(x)|\geq (s-r)(s-1)+r(r-1)=rs+(s-r)^2-s
\end{equation}

\begin{theorem}\label{equal-5-6}
Let $2\leq r<s$ be integers.  
\begin{itemize}
\item
If $n_{bb}(r,s;5)=rs+(s-r)^2+1$, then $r=2$ and $s=3$.
\item
If $n_{bb}(r,s;6)=2(rs+(s-r)^2-s+1)$, then $1.935r < s < 2r-2$.
\end{itemize}
\end{theorem}

\begin{proof}
Let $G$ be an $(r,s;g)$-babi-graph. 
It follows from the proof of Theorem \ref{lowerboundbabigraph} that in the case of equality
$G$ has no thin edge, and $deg_f(x)+deg_f(y)=2(s-r)$ for any fat edge $xy$. 

First, consider the case $g=5$. If the order of $G$ is $rs+(s-r)^2+1$, then equality holds in 
(\ref{tree-5}), so $deg_f(w)=s-r$ for all fat vertex $w$. If $x$ is a fat vertex, then $G_1(x)$
contains $s-r$ fat and $r$ thin vertices. There is no thin edge, hence in $G_2(x)$ there are 
$$(s-r)(s-r-1)+r(r-1)=(s-r)^2+r^2-s$$
fat vertices. So the total number of fat vertices in $G$ is 
$$1+(s-r)+(s-r)^2+r^2-s=(s-r)^2+r^2-r+1.$$
On the other hand, half of the vertices are fat, thus
$$(s-r)^2+r^2-r+1=\frac{rs+(s-r)^2+1}{2}.$$
The solution of this quadratic equation is
$$s=\frac{3r\pm \sqrt{-3r^2+8r-4}}{2}.$$
The discriminant is negative if $r>2.$ If $r=2$, then $s=3$ and $v=8$.

The following construction shows the existence of a $(2,3;5)$-babi-cage of order $8$.
Let $\Gamma$ be the Petersen graph and take an edge $v_1v_2$ of it. Delete the vertices $v_i$ and the edges emanating from $v_i$ for $i=1,2.$ The neighbours of
the $v_i$s are pairwise different vertices, because the girth of $\Gamma $ is 5. 
The obtained graph $\Gamma '$ has 8 vertices, 4 of them (the
remaining neighbours of the $v_i$s) have degree 2, and 4 of them have degree 3.
So $\Gamma $ is a $(2,3;5)$-babi-cage.

\smallskip
Now, consider the $g=6$ case. Let $xy$ be a fat edge of $G$. We may assume without loss of generality
that $1\leq deg_f(x)=s-r-d$ and $s\geq deg_f(y)=s-r+d$ where $d$ is a non-negative integer, $d\leq r$ and $d\leq s-r-1$. Then for any fat neighbour $w$ of $x$ we have 
$$deg_f(w)=2(s-r)-deg_f(x)=s-r+d.$$
There are $s-r-d-1$ fat and $r+d$ thin vertices in $D_1^2(x,y)$. Since $G$ has no thin edge,
$D_2^3(x,y)$ contains
 $$(s-r-1-d)(s-r-1+d)+(r+d)(r-1)=(s-r-1)^2-d^2+(r+d)(r-1)$$
fat vertices. In the same way, we get that there are $s-r+d-1$ fat vertices in $D_1^2(y,x)$ and 
$$(s-r-1)^2-d^2+(r+d)(r-1)$$ 
fat vertices in $D_2^3(y,x)$. So the total number of fat vertices in $G$ is

$$2\left( (s-r)+(s-r-1)^2+r(r-1)-d^2 \right).$$
On the other hand, half of the vertices are fat, thus
$$2\left( (s-r)+(s-r-1)^2+r(r-1)-d^2 \right)=rs+(s-r)^2-s+1.$$
Rearranging this equation, we get
\begin{equation}
\label{diophantine}
s^2-(3r+1)s+4r^2-2d^2+1=0.
\end{equation}

The solution of this quadratic equation is
$$s=\frac{3r+1\pm \sqrt{8d^2-7r^2+6r-3}}{2}.$$
The discriminant is non-negative if 
$$d^2\geq \frac{7}{8}\left( \left( r-\frac{3}{7}\right) ^2+\frac{12}{49}\right) .$$
Thus 
$$d>\sqrt{\frac{7}{8}}\left( r-\frac{3}{7}\right)> 0.935r-0.401 .$$
Since $s-r-1\geq d,$ we get $s>1.935r.$
On the other hand, $d\leq r$ implies 
$$8d^2-7r^2+6r-3\leq r^2+6r-3 < (r+3)^2,$$
hence 
$$s < \frac{3r+1+(r+3)}{2}=2r+2.$$

If $s=2r+h$, then from Equation (\ref{diophantine}) we get
$d^2=r^2+\frac{h-2}{2}r+\frac{h^2-h+1}{2}$. Since the inequalities 
$$r^2-2r+1< r^2+\frac{h-2}{2}r+\frac{h^2-h+1}{2}< r^2$$
hold for $h=-2,-1,0$ and $1$, in these cases $r-1<d<r$, so Equation (\ref{diophantine}) has no integer solution. This proves the second inequality of the statement.
\end{proof}

\end{document}
